% =====================================================================
% The abstract, in English and in Hebrew.
%
% The rules ask for 200 to 500 words in English and 500 to 2000 in
% Hebrew. The Hebrew one is the Introduction without its citations,
% translated and then expanded without adding content.
%
% This file is in UTF-8. Keep it that way.
% =====================================================================

\abstractEnglish{

Neural networks have shown great success in safety-critical
applications, such as medical diagnosis and autonomous driving.
However, even a small change to an input can flip the classification
of the network, and such a change is called an adversarial attack.
Global robustness enables identifying, at
inference time, the inputs that may be susceptible to such attacks.
However, verifying global robustness is a hard task: it reasons about
two copies of the network at once, one for an input and one for the
same input after an attack. This is what has kept global robustness
far less studied in comparison to the extensive research on local robustness,
which reasons about a single input at a time, instead of all possible inputs.

Inspired by the abstractions that scaled local robustness, we propose
lightweight abstractions for scaling global robustness, defined over
two copies of a network. An abstraction describes the relations
between the two copies, and the relations between the different
neurons, which lets the verifier rule out large parts of its
search before the expensive reasoning begins. In addition, we present a
novel conditional-triangle relaxation that extends the popular
triangle approximation to global robustness while reducing its
over-approximation error. We also present, for the first time, an
algorithm for incremental verification of global robustness: a network
is rarely verified once, as it is
retrained and fine-tuned, and every revision raises the same question
again.

We introduce \tool, which enables global robustness verification
on larger networks by combining these complementary abstractions, and
supports incremental verification, leveraging the verification of a
previous network to speed up that of a revised one.
\tool speeds up the exact verifier \baseline by $7.6$
times on average (up to $102.6$ times) when either
terminates, and on the hardest tasks, where both time out, returns
bounds whose gap is $19.9$ times smaller on average (up to
$261.2$ times). When using incremental verification, \tool
speeds up the analysis by $12.5$ times on average (up to
$115.4$ times) when either terminates, and returns
bounds whose gap is $20.9$ times smaller on average (up to $261.2$
times) when both time out.


}

\abstractHebrew{

מסווגים אשר מבוססים על רשתות נוירונים מראים הישגים עם דיוק גבוה מאוד בסיווג הקלטים שהם מקבלים, ועם זאת הם
עדיין פגיעים לדוגמאות אדוורסריאליות: הפרעות קטנות בקלט, אשר משנות את הסיווג של המסווג. כדי לספק ערובות מפני
דוגמאות כאלה, אימות של \emph{חסינות מקומית} מקבל קלט יחיד ונתון, ומאשר שאותו קלט אינו ניתן לסיווג שגוי תחת
אף הפרעה קטנה ההפרעות שעשויות להיות מופעלות עליו. אולם רשת המוטמעת בתחומים הדורשים סף בטיחות גבוה, כגון
נהיגה אוטונומית או הדמיה רפואית, עלולה להתיקל לאחר הטמעתה בקלטים חדשים שמעולם לא ראתה מראש, ועלינו לדעת אילו
מבין כל הקלטים האפשריים ניתנים לסיווג שגוי בשל מתקפה אדוורסריאלית. \emph{חסינות גלובלית} נותנת מענה לצורך
הזה: היא מבטיחה שכל קלט שהמסווג מסווג אותו מעל רמת ביטחון נתונה אינו ניתן להפיכה על ידי שום הפרעה. בזמן
ההסקה, כלומר בזמן שבו המסווג פועל ומסווג קלטים חדשים, כל קלט שהמסווג מסווג אותו מתחת לרמת הביטחון הזו מסומן
כפגיע בפוטנציה, וניתן לטפל בו בהתאם לסימון הזה.

אלא שאימות של חסינות גלובלית הוא משימה קשה. הוא נדרש להסיק על הפרעה של כל קלט אפשרי שהמסווג עשוי לקבל, ולשם
כך עליו לנתח שני עותקים של אותה רשת, האחד עבור הקלט הנקי והאחר עבור ההפרעה שלו, במקום עותק יחיד של הרשת,
כפי שעושים בניתוח של אימות חסינות מקומית. ההכפלה של מספר העותקים של הרשת היא זו שמעלה את הסיבוכיות של הניתוח
בעבור אימות חסינות גלובלית.
לפיכך המאמתים הקיימים ניצבים בפני התפשרות: אימות מדויק של התכונה אורך שעות שלמות אף על רשתות שיש בהן עד
\LRE{2000} נוירונים בלבד, ואילו קירוב-יתר של האימות מאבד דיוק, ולכן הוא מוכיח את החסינות עבור פחות
קלטים.

במחקר זה אנו מבחינים כי ישנן אבסטרקציות קלות משקל, אינטרוולים, זונוטופים ורלקסציות משולש, אשר סייעו להפוך את
אימות החסינות המקומית למשימה ברת-טיפול מבחינה חישובית. הרחבה נאיבית של אבסטרקציות אלה אל החסינות הגלובלית מאבדת דיוק רב
מדי, ושימוש בכל אחת מן האבסטרקציות הללו בפני עצמה, בלי האחרות, מניב האצה מתונה בלבד. רעיון המפתח של עבודה זו
הוא לנצל את הקשר שבין שני העותקים של הרשת, ולעשות זאת באמצעות אבסטרקציות ששומרות יחד על שגיאת קירוב-יתר
קטנה, ובה בעת מאיצות את האימות.

אנו מציגים את \LRE{\tool} \textenglish{(Boosting by Lightweight Encoding of Network Differences),}
כלומר האצה באמצעות קידוד קל משקל של הפרשים בין רשתות, המרחיב את שלוש האבסטרקציות הללו כך שיחולו על שני עותקי
הרשת גם יחד. הראשונה היא אינטרוול הפרש-ההפרעה, החוסם את מידת ההבדל שיכול להיות בין שני העותקים בכל אחד מן
הנוירונים, והוא מרחיב אינטרוולי הפרש אל הפרעות מסוגים שונים, ובכללן טרנספורמציות לינאריות. השנייה היא
זונוטופ, המהדק את חסמי האינטרוול המקובלים המחושבים עבור כל נוירון, והוא עושה זאת באמצעות לכידת התלויות שבין צירופי הערכים
שנוירונים שונים יכולים לקבל יחדיו.
השלישית היא משולש מותנה חדש, המרפה באופן לינארי את אותם נוירונים שמצב ההפעלה שלהם כמעט מוכרע, כלומר נוירונים
שכמעט תמיד פעילים או שכמעט תמיד אינם פעילים, תוך ניצול הקשר שבין שני העותקים. \LRE{\tool} הוא גם
הראשון התומך באימות אינקרמנטלי של חסינות גלובלית: אם רשת שכבר אומתה עוברת אימון מחדש או כיוונון עדין,
\LRE{\tool} מאמת את הרשת המעודכנת על ידי העברת החסמים שהוכחו עבור הרשת הישנה אל הרשת המעודכנת,
ואתחול ראשוני של המשתנים בבעיה החדשה.

אנו מעריכים את \LRE{\tool} באופן ניסויי על רשתות שאומנו על ארבע מערכות נתונים, \LRE{MNIST},
\LRE{\mbox{Fashion-MNIST}}, \LRE{\mbox{CIFAR-10}} ו-\LRE{\mbox{HAR}}, ועל פני שבעה סוגי הפרעות שונים זה
מזה. \LRE{\tool} מאיץ את \LRE{\baseline}, המאמת המדויק הקיים, פי \LRE{7.6} בממוצע,
ועד פי \LRE{102.6}. במשימות הקשות ביותר, שבהן שני המאמתים חורגים ממגבלת הזמן, \LRE{\tool}
מחזיר חסמים הדוקים פי \LRE{19.9} בממוצע מאלה שמחזיר \LRE{\baseline}, ועד פי
\LRE{261.2}. באימות אינקרמנטלי, כלומר כאשר נתונה לו גם תוצאת האימות של הרשת הקודמת,
\LRE{\tool} מאיץ את אימותה של רשת שעברה עדכון פי \LRE{12.5} בממוצע, ועד פי
\LRE{115.4}, ומחזיר חסמים הדוקים פי \LRE{20.9} בממוצע, ועד פי \LRE{261.2}, באותן
משימות שבהן המאמתים חורגים ממגבלת הזמן.

פרק~\ref{ch:background} מגדיר מסווגים המבוססים על רשתות נוירונים, את תכונת החסינות הגלובלית ואת אימותה
המדויק. פרק~\ref{ch:related-work} סוקר את השיטות המאמתות חסינות, ואת המגבלות שבגללן נותרה החסינות הגלובלית
פתוחה. פרק~\ref{ch:key-ideas} מציג את רעיונות המפתח שלנו, ופרק~\ref{ch:blend} מגדיר את הרעיונות הללו באופן
פורמלי. פרק~\ref{ch:evaluation} מעריך את \LRE{\tool} אל מול המאמת המוביל הקודם, הן בזמן הריצה והן
בהידוק החסמים, ומודד את תרומתה של כל אחת מן הטכניקות שלנו בנפרד. פרק~\ref{ch:discussion} מסכם את העבודה
ומציג כיווני המשך שאנו מציעים לעבודה זו.

}
